\documentclass[12pt]{ximera}
\title{Exam 1 (2023)}
\begin{document}
\begin{abstract}
An archived exam on Chapter 1: true/false on the fairness criteria, a Borda count that violates the Majority criterion, Plurality-with-Elimination, and Pairwise Comparisons.
\end{abstract}
\maketitle

\section*{Exam 1 (2023)}

This page combines the two versions of the exam, which differed only in their
true/false statements.

\begin{problem}
True or false?

\begin{prompt}
\textbf{(a)} A candidate who wins a majority of first-place votes will always win under
the Plurality method.
\begin{multipleChoice}
\choice[correct]{True}
\choice{False}
\end{multipleChoice}
\begin{feedback}
True. A majority is more first-place votes than all the others combined.
\end{feedback}
\end{prompt}

\begin{prompt}
\textbf{(b)} A candidate who wins a plurality of first-place votes will always win a
majority.
\begin{multipleChoice}
\choice{True}
\choice[correct]{False}
\end{multipleChoice}
\begin{feedback}
False. With three or more candidates, the most votes can be less than half, as in
Problem 3 below, where A leads with $14$ of $50$.
\end{feedback}
\end{prompt}

\begin{prompt}
\textbf{(c)} A candidate who wins a majority of first-place votes is always a Condorcet
candidate.
\begin{multipleChoice}
\choice[correct]{True}
\choice{False}
\end{multipleChoice}
\begin{feedback}
True. Those same majority voters prefer the candidate over every opponent.
\end{feedback}
\end{prompt}

\begin{prompt}
\textbf{(d)} A Condorcet candidate always wins under the Pairwise Comparisons method.
\begin{multipleChoice}
\choice[correct]{True}
\choice{False}
\end{multipleChoice}
\begin{feedback}
True. Winning every comparison gives the most points.
\end{feedback}
\end{prompt}

\begin{prompt}
\textbf{(e)} A candidate who wins a majority of first-place votes always wins under the
Borda count method.
\begin{multipleChoice}
\choice{True}
\choice[correct]{False}
\end{multipleChoice}
\begin{feedback}
False. Problem 2 below is a counterexample: H has $21$ of $40$ first-place votes but
loses the Borda count.
\end{feedback}
\end{prompt}

\begin{prompt}
\textbf{(f)} There must be a Condorcet candidate in every election.
\begin{multipleChoice}
\choice{True}
\choice[correct]{False}
\end{multipleChoice}
\begin{feedback}
False. Head-to-head results can form a cycle, with A beating B, B beating C, and C
beating A.
\end{feedback}
\end{prompt}

\begin{prompt}
\textbf{(g)} A Condorcet candidate always wins under the Plurality-with-Elimination
method.
\begin{multipleChoice}
\choice{True}
\choice[correct]{False}
\end{multipleChoice}
\begin{feedback}
False. A Condorcet candidate with few first-place votes can be eliminated early.
\end{feedback}
\end{prompt}
\end{problem}

\begin{problem}
Three candidates, H, D, and C, are running for mayor.
\begin{center}
\begin{tabular}{c|ccc}
Number of voters & 21 & 10 & 9\\\hline
1st & H & D & D\\
2nd & D & C & H\\
3rd & C & H & C
\end{tabular}
\end{center}
Use the (standard) Borda count to find the points and the winner.

H: $\answer{91}$ \quad D: $\answer{99}$ \quad C: $\answer{50}$ \qquad Winner: $\answer{D}$

\begin{feedback}
\[
\begin{aligned}
\text{H:}&\ 21\cdot3+10\cdot1+9\cdot2=91\\
\text{D:}&\ 21\cdot2+10\cdot3+9\cdot3=99\\
\text{C:}&\ 21\cdot1+10\cdot2+9\cdot1=50
\end{aligned}
\]
D wins, even though H has a majority of first-place votes. This violates the Majority
criterion.
\end{feedback}
\end{problem}

\begin{problem}
Campers at a zoo's summer camp rank their favorite animals: alligators (A), baboons (B),
camels (C), and dolphins (D).
\begin{center}
\begin{tabular}{c|ccccc}
Number of voters & 14 & 12 & 6 & 5 & 13\\\hline
1st & A & D & B & B & C\\
2nd & B & A & C & D & A\\
3rd & C & B & A & C & D\\
4th & D & C & D & A & B
\end{tabular}
\end{center}
Determine the winner under Plurality-with-Elimination.

Eliminated in round 1: $\answer{B}$ \quad in round 2: $\answer{A}$ \qquad Winner:
$\answer{C}$

\begin{feedback}
A majority of $50$ is $26$.
\begin{itemize}
\item Round 1: A $14$, C $13$, D $12$, B $11$. Eliminate B: $6$ votes go to C and $5$
      to D.
\item Round 2: C $19$, D $17$, A $14$. Eliminate A; its $14$ voters rank C next.
\item Round 3: C $33$, D $17$. Camels win.
\end{itemize}
\end{feedback}
\end{problem}

\begin{problem}
A soccer team returning from an away game is choosing where to stop for dinner: Wendy's
(W), Taco Bell (T), or Subway (S).
\begin{center}
\begin{tabular}{c|cccc}
Number of voters & 10 & 8 & 5 & 7\\\hline
1st & W & S & T & T\\
2nd & T & T & S & W\\
3rd & S & W & W & S
\end{tabular}
\end{center}
Use Pairwise Comparisons to determine where the team stops.

$\answer{T}$

\begin{feedback}
T beats W $20$--$10$ and S $22$--$8$; W beats S $17$--$13$. Taco Bell wins with $2$
points.
\end{feedback}
\end{problem}

\end{document}
